% Part: lambda-calculus
% Chapter: introduction
% Section: reduction

\documentclass[../../../include/open-logic-section]{subfiles}

\begin{document}

\olfileid{lam}{int}{red}
\olsection{کاهش ترم‌های لامبدا}

با ترم‌های لامبدا چه می‌توان کرد؟ می‌توان آن‌ها را ساده کرد. اگر~$M$
و~$N$ دو ترم لامبدای دلخواه و~$x$ متغیری دلخواه باشند، می‌توانیم
$\Subst{M}{N}{x}$ را برای نمایش نتیجهٔ جایگزینی~$N$ به‌جای~$x$ در~$M$
به کار ببریم، البته پس از تغییر نام هر متغیر مقید~$M$ که بعد از جانشانی
با متغیرهای آزاد~$N$ تداخل پیدا کند. برای نمونه،
\[
\Subst{(\lambd[w][xxw])}{yyz}{x} = \lambd[w][(yyz)(yyz)w].
\]

\begin{digress}
نمادگذاری‌های دیگر برای جانشانی عبارت‌اند از~$[N/x]M$، $[x/N]M$ و نیز
$M[x/N]$. مراقب باشید!
\end{digress}

به‌طور شهودی، $(\lambd[x][M])N$ و~$\Subst{M}{N}{x}$ معنای یکسانی
دارند؛ عمل جایگزین‌کردن ترم نخست با ترم دوم \emph{انقباض~$\beta$}
نامیده می‌شود. $(\lambd[x][M])N$ را \emph{ردکس} و
$\Subst{M}{N}{x}$ را \emph{منقبضۀ} آن می‌نامند. به‌طور کلی، اگر بتوان
ترمی چون~$P$ را به~$P'$ با انقباض~$\beta$ یک زیرترم تغییر داد، می‌گوییم
$P$ \emph{با یک گام کاهش~$\beta$ به~$P'$ کاهش می‌یابد} و می‌نویسیم
$P \redone P'$. اگر بتوان از~$P$ با تعدادی کاهش یک‌گامی (شاید صفر گام)
به~$P'$ رسید، آنگاه می‌گوییم~$P$ \emph{به روش~$\beta$ به~$P'$ کاهش
می‌یابد}؛ به نماد،
$P \red P'$. ترمی که دیگر نتوان آن را~$\beta$-کاهش داد
\emph{$\beta$-کاهش‌ناپذیر} یا \emph{$\beta$-نرمال} نامیده می‌شود. هرگاه
بافت روشن باشد، به‌جای «$\beta$-کاهش می‌یابد» و مانند آن، فقط می‌گوییم
«کاهش می‌یابد».

چند نمونه را در نظر بگیریم.
\begin{enumerate}
\item داریم
\begin{align*}
(\lambd[x][xxy]) \lambd[z][z] & \redone (\lambd[z][z])(\lambd[z][z]) y \\
& \redone (\lambd[z][z]) y \\
& \redone y.
\end{align*}
\item «ساده‌کردن» یک ترم می‌تواند آن را پیچیده‌تر کند:
\begin{align*}
(\lambd[x][xxy])(\lambd[x][xxy]) & \redone (\lambd[x][xxy])(\lambd[x][xxy])y \\
& \redone (\lambd[x][xxy])(\lambd[x][xxy])yy \\
& \redone \dots
\end{align*}
\item همچنین ممکن است ترم را بی‌تغییر بگذارد:
\[
(\lambd[x][xx])(\lambd[x][xx]) \redone (\lambd[x][xx])(\lambd[x][xx]).
\]
\item افزون بر این، برخی ترم‌ها را می‌توان به بیش از یک شیوه کاهش داد؛
  برای نمونه،
\[
(\lambd[x][(\lambd[y][yx]) z)] v \redone (\lambd[y][yv]) z
\]
با منقبض‌کردن بیرونی‌ترین کاربرد؛ و
\[
(\lambd[x][(\lambd[y][yx]) z)] v \redone (\lambd[x][zx]) v
\]
با منقبض‌کردن درونی‌ترین کاربرد. بااین‌حال، توجه کنید که در این حالت هر
دو ترم در ادامه به ترم یکسان~$zv$ کاهش می‌یابند.
\end{enumerate}

نتیجهٔ نهایی در نمونهٔ آخر تصادفی نیست؛ بلکه ویژگی عمیق و مهمی از حساب
لامبدا را نشان می‌دهد که «ویژگی چرچ--راسر» نام دارد.

\end{document}
